\section*{Overview}

The Mobius function \(\mu(n)\) is the standard tool for removing overcounting across divisors. In practice it appears
when a problem wants ``gcd equals \(1\)'' or ``exactly this divisor'' and the easier count is over multiples instead.

\section*{When to Use It}

Use it when:

\begin{itemize}
  \item divisor inclusion-exclusion is visible,
  \item you can count objects divisible by \(d\) more easily than objects with gcd exactly \(1\),
  \item a multiplicative-function sieve is already nearby.
\end{itemize}

\section*{Core Idea}

The Mobius function is:

\begin{itemize}
  \item \(0\) if \(n\) has a squared prime factor,
  \item \(1\) if \(n\) is a product of an even number of distinct primes,
  \item \(-1\) if \(n\) is a product of an odd number of distinct primes.
\end{itemize}

Its main identity is
\[
\sum_{d \mid n} \mu(d) =
\begin{cases}
1 & n = 1, \\
0 & n > 1.
\end{cases}
\]

\section*{Key Insight}

That identity is a switch that isolates gcd \(=1\). When you sum over divisors with coefficient \(\mu(d)\), every number
with a larger common divisor cancels out.

\section*{Worked Problem}

\paragraph{Problem.}
How many pairs \((x, y)\) with \(1 \le x \le n\), \(1 \le y \le m\) satisfy \(\gcd(x, y) = 1\)?

\paragraph{Why Mobius fits.}
For one divisor \(d\), the number of pairs where \(d\) divides both numbers is
\[
\left\lfloor \frac{n}{d} \right\rfloor \left\lfloor \frac{m}{d} \right\rfloor.
\]
Applying Mobius inversion gives
\[
\sum_{d=1}^{\min(n,m)} \mu(d)\left\lfloor \frac{n}{d} \right\rfloor \left\lfloor \frac{m}{d} \right\rfloor.
\]

\section*{Correctness Intuition}

Each pair contributes to all common divisors of \(\gcd(x,y)\). The Mobius coefficients cancel every pair whose gcd is
greater than \(1\), leaving only the coprime pairs.

\section*{Complexity Analysis}

Computing \(\mu\) up to \(N\) by linear sieve takes \(O(N)\). The direct summation formula above is \(O(N)\), with
further harmonic-speed tricks possible in heavier problems.

\section*{Implementation}

The code computes \(\mu\) with a linear sieve and includes the coprime-pair formula.

\section*{Common Pitfalls}

\begin{itemize}
  \item Mixing Mobius inversion with Euler phi; they answer different questions.
  \item Forgetting that squareful numbers contribute \(0\).
  \item Expanding divisor sums naively without checking whether a harmonic grouping is needed.
\end{itemize}

\section*{Variants / Extensions}

\begin{itemize}
  \item Mobius inversion on Dirichlet convolution.
  \item Counting tuples with gcd \(1\).
  \item Summatory multiplicative-function problems with prefix tricks.
\end{itemize}

\section*{Practice Problems}

\begin{itemize}
  \item Count coprime pairs.
  \item Count pairs with gcd exactly \(g\).
  \item Any divisor-sum problem where cancellation over multiples is the main obstacle.
\end{itemize}

\section*{References}

\begin{itemize}
  \item \href{https://cp-algorithms.com/algebra/mobius-function.html}{cp-algorithms: Mobius Function}.
  \item \href{https://oiwiki.com/math/number-theory/mobius/}{OI Wiki: Mobius Function}.
  \item \href{https://codeforces.com/catalog}{Codeforces Catalog}.
\end{itemize}
